\documentclass[11pt]{letter}
\usepackage[utf8]{inputenc}
\usepackage{geometry}
\geometry{a4paper, margin=1in}
\usepackage{hyperref}

\address{Florian Rzepka\\
Technical University of Berlin\\
Chair of Electrical Energy Storage Technology\\
Einsteinufer 11\\
10587 Berlin, Germany\\
\href{mailto:florian.rzepka@tu-berlin.de}{florian.rzepka@tu-berlin.de}}

\date{\today}

\begin{document}
\begin{letter}{Editor-in-Chief\\
Engineering Applications of Artificial Intelligence\\
Elsevier}

\opening{Dear Editor-in-Chief,}

Please find enclosed our manuscript entitled \textbf{``Performance Benchmarking of Embedded Battery State Estimation: A Measurement-Only Comparison of Direct-Measurement, Hybrid, and Data-Driven Estimators''}, which we submit for exclusive consideration in \textit{Engineering Applications of Artificial Intelligence}.

The manuscript addresses a practical gap in embedded battery management: estimator selection is often discussed through clean-data accuracy alone, whereas deployment quality depends equally on behavior under realistic measurement disturbances. To address this, the paper introduces a controlled measurement-only benchmark that compares four embedded-ready estimator classes under the same full-cell trajectory, the same disturbance definitions, and the same online execution assumptions. The study evaluates direct-measurement, hybrid direct-measurement, hybrid equivalent-circuit, and data-driven estimators under ADC quantization, current bias, additive noise, missing samples, irregular sampling, burst dropout, voltage spikes, and initialization mismatch.

The main contribution is a cross-class performance benchmark that separates baseline accuracy, disturbed-operation robustness, and recovery behavior under physically interpretable sensor-side faults. Rather than promoting a single universal winner, the results expose complementary strengths: the hybrid direct-measurement model leads under clean conditions, the hybrid equivalent-circuit model is the strongest overall robustness choice, and the data-driven model is particularly attractive for missing-sample resilience and strict-band recovery after initialization mismatch. This makes the paper directly relevant to the readership of \textit{EAAI}, because it connects data-driven and hybrid AI-enabled estimation with deployment-oriented evaluation and decision support for embedded energy systems.

We confirm that this manuscript has not been published elsewhere and is not under consideration by another journal. All authors have approved the manuscript and agree with its submission.

Thank you for your consideration.

\closing{Sincerely,}

Florian Rzepka\\
Technical University of Berlin
\end{letter}
\end{document}
